% !TeX program = pdflatex
\documentclass[11pt,a4paper]{article}

% --- Język i czcionki ---
\usepackage[T1]{fontenc}
\usepackage[utf8]{inputenc}
\usepackage[polish]{babel}
\usepackage{lmodern}
\usepackage[final]{microtype}

% --- Układ strony i nagłówki/stopki ---
\usepackage[a4paper,margin=2.5cm]{geometry}
\usepackage{fancyhdr}
\pagestyle{fancy}
\fancyhf{} % czyść wszystko
\lhead{Etap 2.}
\rhead{\leftmark}
\cfoot{\thepage}

% --- Matematyka ---
\usepackage{amsmath,amssymb,amsthm,mathtools}
\usepackage{siunitx} % jeżeli potrzebne jednostki
\sisetup{locale = PL} % separator 1,23 (polski), można zmienić na locale=PL gdy dostępne
\usepackage{bm}       % pogrubione symbole

% --- Odsyłacze i tytuły ---
\usepackage[hidelinks]{hyperref}
\usepackage[nameinlink,capitalise,noabbrev]{cleveref}
\usepackage{csquotes}

% --- Listy i drobiazgi ---
\usepackage{enumitem}
\setlist{noitemsep,topsep=3pt}
\usepackage{xcolor}

% --- Numeracja i wygląd sekcji ---
\numberwithin{equation}{section}
\setcounter{secnumdepth}{3}
\setcounter{tocdepth}{2}

% --- Meta ---
\title{Sprawozdanie z etapu 2.}
\author{Jakub Kogut}
\date{\today}

\begin{document}
\maketitle
\section{Wymagania Funkcjonalne}
Nasz projekt zakłada stworzenie aplikacji mobilnej dla systemu Android z natępującymi funkcjonalnościami:
\begin{itemize}
    \item Rejestracja i logowanie użytkowników.
    \item Modyfikacja profilu użytkownika (wybranie awatatu).
    \item System poziomów i odznak za postępy w nauce.
    \item Przeglądanie swoich postępów w nauce (poziomy, zdobyte odznaki itp.).
    \item Nauka z kursów, podzielonych w następujący sposób:
        \[
        \text{Kurs} \rightarrow \text{Lekcje}\quad \text{Lekcja} \rightarrow \text{Slaidy} + \text{Quiz na koniec lekcji}
        \]
    \item Quizy będą zawierały pytania losowane z puli, otwarte (wartości liczbowe) oraz zamknięte.
    \item Ekran końcowy z wynikami podsumowującymi wyniki quizu.
    \item Forum dyskusyjne dla użytkowników podzielone na kategorie z systemem komentarzy i odpowiedzi, ocen postów.
    \item Zakładka z quizami do samodzielnej nauki.
\end{itemize}
\section{Podział na zadania}
Planujemy podzielić projekt na dwa obszary:
\begin{itemize}
    \item Backend (Django + DRF):
        \begin{enumerate}
            \item Analiza i planowanie bazy danych.
            \item Zaplanowanie wymaganych endpointów API.
            \item Implementacja modeli bazy danych (podział na sekcje: Kursy, Quizy, Użytkownicy, Forum).
            \item Implementacja endpointów API (zgodnie z w.w. podziałem)
            \item Stworzenie przyjaznego dla administratora panelu administracyjnego (dodawanie treści).
            \item Testowanie i dokumentacja API.
            \item Wystawienie API na serwerze.
        \end{enumerate}
    \item Frontend (Aplikacja mobilna):
        \begin{enumerate}
            \item Stworzenie logiki odpowiedzialnej za komunikację z API.
            \item Tworzenie przyjaznego interfejsu użytkownika (UI/UX).
                \begin{itemize}
                    \item Ekran startowy,
                    \item Ekran rejestracji/logowania,
                    \item Ekran profilu użytkownika,
                    \item System nauki: ekrany kursów, lekcji, slaidów, quizów i przepływ między nimi,
                    \item Ekran wyników quizu,
                    \item Ekran forum dyskusyjnego,
                    \item Ekran z quizami do samodzielnej nauki.
                \end{itemize}
            \item Integracja z backendem (wykorzystanie endpointów API).
            \item Obsługa błędów i wyjątków (np. brak połączenia z internetem).
            \item Testowanie aplikacji na różnych urządzeniach.
        \end{enumerate}
\end{itemize}
Do każdego z obszarów przydzieliśmy po dwóch członków zespołu, którzy będą odpowiedzialni za jego realizację:
\begin{itemize}
    \item Backend:
        \begin{itemize}
            \item Jakub Kogut
            \item Wiktor Koczkodaj
        \end{itemize}
    \item Aplikacja mobilna:
        \begin{itemize}
            \item Szymon Hładyszewski
            \item Wojciech Typer
        \end{itemize}
\end{itemize}
\section{Harmonogram Prac}
Planowany harmonogram prac przedstawia się następująco:
%dolacz gantt.png
\begin{figure}[h]
    \centering
    \includegraphics[width=\textwidth]{gantt.png}
    \caption{Harmonogram prac nad projektem}
    \label{fig:gantt}
\end{figure}
\section{Wybrane Technologie}
Do wykonania backendu planujemy użyć frameworka Django wraz z Django REST Framework (DRF) do tworzenia API. Baza danych zostanie zaimplementowana przy użyciu PostgreSQL, będzie zarządzana przez Django. Framework ten w łatwy sposób pozwala na tworzenie dobrej dokumentacji API, panelu administratora oraz jest wysoce modularny i skalowalny. 
\\
Natomiast do stworzenia samej aplikacji mobilnej planujemy użyć:
\begin{itemize}
    \item \textbf{Kotlin}: nowoczesny język programowania stworzony z myślą o tworzeniu aplikacji mobilnych. Wyróżnia się duża ilością popularnych frameworków, dobrze obsługuje wielowątkowość oraz jest mocno pokrewny z Javą.
    \item \textbf{Jetpack Compose}: framework UI, który umożliwia:
        \begin{itemize}
            \item tworzenie UI bez użycia XML,
            \item pisanie całej logiki UI w Kotlinie,
            \item intuicyjne zarządzanie stanem aplikacji,
            \item łatwą integrację z innymi narzędziami Androida.
        \end{itemize}
\end{itemize}
            
\end{document}

